\documentclass[withoutpreface,bwprint]{cumcmthesis}
\usepackage[framemethod=TikZ]{mdframed}
\usepackage{amstext} 
\usepackage{amssymb}
\usepackage{amsmath}
\usepackage{booktabs}
\usepackage{bm}
\usepackage{url}
\usepackage{subcaption}
\usepackage{mathtools} 
\usepackage[linesnumbered,ruled,lined]{algorithm2e}
\usepackage{makecell}
\usepackage{longtable}
\usepackage{metalogo}
\usepackage{setspace}

\makeatletter
\providecommand\IfPDFManagementActiveF[1]{#1}
\makeatother

\usepackage{caption}
\usepackage{multirow}
\usepackage{float}
\usepackage{graphicx}
\usepackage[linesnumbered,ruled,lined]{algorithm2e}
\usepackage{threeparttable}
\usepackage[square,sort,comma,numbers]{natbib}
\usepackage{diagbox}
\usepackage{color}
\usepackage{transparent}
\usepackage{import}
\usepackage{hyperref}





%\special{dvipdfmx:config z 0} %取消PDF压缩，加快速度，最终版本生成的时候最好把这句话注释掉
\title{基于}
\begin{document}
\maketitle
\thispagestyle{empty}
\begin{abstract}\setlength{\parskip}{0.3\baselineskip}



\textbf{问题1}

\textbf{问题2}

\textbf{问题3}

\textbf{本文优点：}
    \keywords{ \textbf{}  \quad  \textbf{ }  \quad \textbf{} \quad \textbf{}\quad \textbf{}}
\end{abstract}
\thispagestyle{empty}
\tableofcontents
\thispagestyle{empty}
\newpage
\setcounter{page}{1} 




\section{问题重述}
\subsection{问题背景}



\subsection{问题概述}


\section{模型假设}





\section{符号说明}
\begin{table}[!h]
    \centering
    \begin{tabular}{p{2.5cm}<{\centering}p{12cm}<{\centering}p{2.5cm}<{\centering}}
    \toprule[1.5pt]
    符号 & 含义\\
    \midrule[1pt]
    $w_1 $ &  不满意程度的权重    \\
    $w_2 $ &    负面影响程度的权重   \\
   
    \bottomrule[1.5pt]
    \end{tabular}
\end{table}
注：表中未说明的符号以首次出现处为准


\section{问题一模型的建立与求解}

\subsection{问题分析}

\subsection{模型建立}

\subsection {模型求解和结果分析}

\subsection{模型检验}




%此部分用于预留常用代码
%%%%%%%%%%%%%%%%%%%%%%%%%%%%%%%%%
%          伪代码
%%%%%%%%%%%%%%%%%%%%%%%%%%%%%%%%%
\iffalse
\begin{algorithm}
		\caption{Function ImperialistCompetitveAlgorithm(x)}
		\LinesNumbered
		\textbf{Initialize}
		\tcc{初始化：设置多个初始国家解$X_n,n=1,...,n_max$。从中选取$N$个最好质量的解使其作为$N$个帝国的殖民国家，$X_1,X_2,...,X_N$, 为每个殖民国家分配殖民国家势力比例数量的殖民地得到$N$个帝国$E_1,E_2,...,E_N$}
		iteration = 1;
		\textbf{Repeat}
		\For{$i \in [1,N]$}{
			EmpireAssimilate($E_i$) \tcc{帝国同化}
			ColoniesRevolve($E_i$) \tcc{殖民地革命}
			EmpirePosses($E_i$) \tcc{殖民国家替换}
			CaculateNewCost($E_i$) \tcc{重新计算国家权力}
		}
		EmpireUnite({$E_1, E_2, \dots, E_N$}) \tcc{合并相似帝国}
		EmpireCompetition({$E_1,E_2,...,E_N$}) \tcc{帝国竞争，分配殖民地}
		iteration $\leftarrow$ iteration$+1$; {\mbox{进入下一次迭代}} \newline
		until \mbox{只有一个帝国剩下/达到迭代次数。}
	\end{algorithm}
\fi

%%%%%%%%%%%%%%%%%%%%%%%%%%%%%%%%%
%          插图
%%%%%%%%%%%%%%%%%%%%%%%%%%%%%%%%%
\iffalse
%         单幅图
\begin{figure}[H]
    \centering
    \includegraphics[width=0.7\textwidth]{}
    \caption{}
    \label{fig:}
\end{figure}
\fi
\iffalse
%         两图并列
\begin{figure}[!h]
	\centering
	\begin{minipage}[c]{0.48\textwidth}
    	\centering
	    \includegraphics[width=1\textwidth]{}
	    \caption{}
	    \label{fig:}
    \end{minipage}
    \begin{minipage}[c]{0.48\textwidth}
		\centering
        \includegraphics[width=1\textwidth]{}
        \caption{}
        \label{fig:}
	\end{minipage}
\end{figure}
\fi
\iffalse
%       四幅子图
\begin{figure}[H]
	\centering
	\subfigure[]{\includegraphics[width=.43\textwidth]{}}
	\subfigure[]{\includegraphics[width=.43\textwidth]{}}\\
	\subfigure[]{\includegraphics[width=.43\textwidth]{}}
	\subfigure[]{\includegraphics[width=.43\textwidth]{}}
	\caption{}
	\label{fig:}
\end{figure}
\fi

%%%%%%%%%%%%%%%%%%%%%%%%%%%%%%%%%
%           表格
%%%%%%%%%%%%%%%%%%%%%%%%%%%%%%%%%
\iffalse
\begin{table}[!h]
	\caption{}
	\label{tab:}
	\centering
	\begin{tabular}{p{3cm}<{\centering}p{3cm}<{\centering}p{3cm}<{\centering}}
		\toprule[1.5pt]
		   &     &   \\ 
		\midrule[1pt]
		   &     &    \\
		\bottomrule[1.5pt]
	\end{tabular}
\end{table}
\fi
%%%%%%%%%%%%%%%%%%%%%%%%%%%%%%%%%
%           公式
%%%%%%%%%%%%%%%%%%%%%%%%%%%%%%%%%
\iffalse
\begin{equation}
    \left\{
         \begin{array}{lr}
       
       \frac{R}{\sin\alpha_1}=\frac{r}{\sin(\alpha_1+\varphi)}, \\
       \frac{R}{\sin\alpha_2}=\frac{r}{\sin(\alpha_2+\left|\theta -\varphi\right|)}.\\
 \end{array}
  \right.
 \end{equation}


\fi
\iffalse
\begin{aligned}
	\max & \sum_{j} X_{12 j} \times 1000+10000-\sum_{k}\left[\left(T_{k 1}+\sum_{j=2}^{27} 2 T_{k} j\right) a_{k}\right], \\
	\text { s.t. } & \left\{\begin{array}{l}
	\sum_{j=0}^{30} x_{27 j}=1, \\
	\sum_{j=0}^{30} x_{1 j} \geq 1, \\
	x_{i j} x_{i^{\prime} j+1} \leq d_{i i^{\prime}}, \\
	V_{j} \leq 2+M x_{12 j}, \\
	Q_{k j}+T_{k j}-S_{k j} \geq 0, \\
	Q_{k j+1}=Q_{k j}+T_{k j}-S_{k j}, \\
	\\
	3 Q_{1 j}+2 Q_{2 j} \leq 1200 .
	\end{array}\right.
	\end{aligned}
\fi






\section{问题二模型的建立与求解}

\subsection {模型建立}

\subsection {模型求解}




\section{问题三模型的建立与求解}


\subsection {模型建立}


\subsection {模型求解}



\subsection {结果分析}


\section{模型评价}
\subsection {模型优点}


\subsection {模型缺点}


\subsection{模型评价}






\newpage
\begin{thebibliography}{99}
    \bibitem{1}
    \bibitem{2}
    \bibitem{3}
    \bibitem{4}
    \bibitem{5}
    \bibitem{7}
    \bibitem{8}
    \bibitem{9}
    \bibitem{10}
    \bibitem{11}
    \bibitem{12}
    \bibitem{13}
    \bibitem{14}
        
\end{thebibliography}



\newpage
\begin{appendices}


\section{问题一源代码}
\begin{lstlisting}[language=lingo]


\end{lstlisting}


\section{问题二源代码}
\begin{lstlisting}[language=lingo]

  
\end{lstlisting}

\section{问题三源代码}
\begin{lstlisting}[language=lingo]

\end{lstlisting}

\begin{lstlisting}[language=python]



\end{lstlisting}
\end{appendices}
\end{document}